\documentclass[10pt,a4paper]{amsart}
\usepackage[margin=19mm]{geometry}
\usepackage{amsmath,amssymb,mathtools}
\usepackage[scr=rsfs,cal=euler]{mathalfa}
\usepackage{xfrac}
\usepackage[unicode,hidelinks,bookmarks=false]{hyperref}
\newcommand{\ie}{i.e.\ }
\newcommand{\cf}{cf.\ }
\newcommand{\resp}{respectively\ }
\newcommand{\Subsection}{\subsection}
\providecommand{\nobreakdash}{\nobreak\hbox{-}}
\setlength{\emergencystretch}{2em}
\multlinegap0pt
\usepackage{amssymb}
\usepackage{mathtools}
\usepackage[normalem]{ulem}
\usepackage{xcolor}
\usepackage{tikz}
\usepackage{cancel}
\usepackage[shortlabels]{enumitem}
\usepackage{float}
\newtheorem{lemma}{Lemma}
\newtheorem{theorem}{Theorem}
\newcommand\Set[2]{\left\lbrace #1 \mid #2 \right\rbrace}
\newcommand{\calS}{\mathcal{S}}
\newcommand\set[1]{\left\lbrace #1 \right\rbrace}
\title{Open Neighborhood Ideals of Well-Totally Dominated Trees are Cohen-Macaulay}
\author{Jounglag Lim, James Gossell, Keri Ann Sather-Wagstaff}
\date{Source: arXiv:2510.19707v4 (CC BY 4.0)}
\begin{document}
\maketitle

\noindent\emph{Source and licence.} Open Neighborhood Ideals of Well-Totally Dominated Trees are Cohen-Macaulay. Version arXiv:2510.19707v4 (CC BY 4.0), distributed by arXiv under the Creative Commons Attribution 4.0 International licence, http://creativecommons.org/licenses/by/4.0/, as recorded in the arXiv licence metadata for this exact version and shown on the abstract page https://arxiv.org/abs/2510.19707v4. Full licence notice: \texttt{licenses/arxiv-2510.19707-CC-BY-4.0.txt}.\par\smallskip

\noindent\textit{Editorial scope.} Counted unit: the article's theorem shellable (source0.tex lines 984--987). The article's complete original proof of it is delivered with it (source0.tex lines 989--1044, one \texttt{proof} environment of 56 lines). No mathematics is written, repaired or re-derived: the statement and the proof are the article's own lines, in the article's own notation, and no retained line is substituted or re-flowed.  Retained uncounted as the source-local context the proof needs: lines 839--846; lines 916--923; lines 948--958.  Nothing was omitted from any retained span, so the package carries no omission record.  No retained line contains a citation, so the wrapper carries no bibliography and no bibliography entry.  No retained line contains a figure environment, an image-inclusion command or drawing code, so no image is delivered and the package declares no asset dependency.  Editorial formatting only, in addition: an AMS \texttt{amsart} preamble that restates the article's own declarations and macros used by the retained lines, namely the environments \texttt{lemma}, \texttt{theorem} on separate counters, together with the standard packages those lines need.  The retained environments print in this document as Lemma 1, Theorem 1 in order of appearance, whereas the article prints the same items with the numbering of its own full document.  The complete native source of the article as distributed in the arXiv e-print for this exact version is bundled unchanged as \texttt{sources/187591/source0.tex}.
\begin{lemma}\label{lem. |N(s) cap D| = 1}
    Let $T$ be an WTD balanced tree of height 3, and let $D$ be a minimal odd-TDS.
    Let $s \in V$ be a support vertex.
    Then $|N(s) \cap D| = 1$.
    Therefore, every minimal odd-TDS of $T$ is of the form
    $\set{u_{1,a_1},u_{2,a_2},\dots,u_{p,a_{p}}}$
where $1 \leq a_i \leq k_i$.
\end{lemma}

\begin{lemma}\label{lem. F1 cap F2 size}
    Let $T$ be an WTD balanced tree of height 3.
    Let $F_1, F_2 \in \mathcal{F}(\calS_{even}(T))$ with $\nu_T(F_1) = (a_1,\dots,a_p)$ and $\nu_T(F_2) = (b_1,\dots,b_p)$.
    Then
    $$
    |F_1 \cap F_2| = \dim(\calS_{even}(T)) + 1 - |\Set{i}{a_i \neq b_i, 1\leq i \leq p}|.
    $$
\end{lemma}

\begin{lemma}\label{lem. non-zero to 1 is TDS}
    Let $T$ be an WTD balanced tree of height 3, and let 
    $$
    D = \set{u_{1,a_1},\dots,u_{p,a_p}} \subseteq V
    $$ 
    be a minimal odd-TDS in $T$.
    Suppose that there exists some index $i$ such that $a_i \neq 1$.
    Then for any $c \in [k_i]$, the set
    $D' = (D\setminus \set{u_{i,a_i}}) \cup \set{u_{i,c}}$
    is also a minimal odd-TDS.
\end{lemma}

\begin{theorem}\label{thm. even stable complex is shellable}
    Let $T$ be an WTD balanced tree of height 3.
    Then the total order $<_T$ on $U_T$ restricted to $\mathcal{F}(\calS_{even}(T))$ is a shelling of $\calS_{even}(T)$.
\end{theorem}

\begin{proof}
    Since $\mathcal{F}(\calS_{even}(T)) \subseteq U_T $ by Lemma~\ref{lem. |N(s) cap D| = 1}, $<_T$ is a total order on $\mathcal{F}(\calS_{even}(T))$.
    Now index the facets in $\calS_{even}(T)$ under $<_T$ and write
    $$
    \mathcal{F}(\calS_{even}(T)) = \set{F_1,F_2,\dots,F_m}.
    $$
    To show that $<_T$ is a shelling, suppose that there exist indicies $i,j$ with $1 \leq i < j \leq m$ such that $|F_i \cap F_j| < \dim(\calS_{even}(T))$.
    We need to find a facet $F \in \mathcal{F}(\calS_{even}(T))$ such that $F <_T F_j$, $F_i \cap F_j \subseteq F \cap F_j$, and $|F \cap F_j| = \dim(\calS_{even}(T))$.
    Write $\nu_T(F_i) = (a_1,\dots,a_p)$ and $\nu_T(F_j) = (b_1,\dots,b_p)$.
    Since $F_i <_T F_j$, $\nu_T(F_i) <_P \nu_T(F_j)$, so there are two cases to consider.\newline
    
    \noindent\emph{Case 1}: $\ell(\nu_T(F_i)) > \ell(\nu_T(F_j))$.
    Then $\nu_T(F_i)$ has more 1's then $\nu_T(F_j)$.
    Hence there exists some index $k$ such that $a_k = 1$ and $b_k \neq 1$.
    Now consider 
    $$
    F = V_{even}\setminus \set{u_{1,b_1},\dots,u_{k-1,b_{k-1}}, u_{k,1}, u_{k+1,b_{k+1}},\dots,u_{p,b_p}} \in U_T 
    $$
    with
    $$
    \nu_T(F) = (b_1,\dots,b_{k-1},1,b_{k+1},\dots,b_p).
    $$
    Then $F \in \mathcal{F}(\calS_{even}(T))$ by Lemma~\ref{lem. non-zero to 1 is TDS}.
    Since $\ell(\nu_T(F)) = \ell(\nu_T(F_j)) + 1$, we get $F <_T F_j$.
    Since we have
    $$
    F_i \cap F_j = V_{even} \setminus \left( \set{u_{1,a_1},\dots,u_{p,a_p}} \cup \set{u_{1,b_1},\dots,u_{p,b_p}}\right)
    $$
    and
    \begin{align*}
        F \cap F_j &= V_{evem} \setminus \left(\set{u_{1,b_1},\dots,u_{k-1,b_{k-1}}, u_{k,1}, u_{k+1,b_{k+1}},\dots,u_{p,b_p}} \cup \set{u_{1,b_1},\dots,u_{p,b_p}} \right)\\ 
        &= V_{even} \setminus \left( \set{u_{k,1}}\cup\set{u_{1,b_1},\dots,u_{p,b_p}}\right)
    \end{align*}
    the choice of $k$ giving $u_{k,1} = u_{k,a_k} \in \set{u_{1,a_1},\dots,u_{p,a_p}}$ implies 
    $$
    \left( \set{u_{k,1}}\cup\set{u_{1,b_1},\dots,u_{p,b_p}}\right) \subseteq \left( \set{u_{1,a_1},\dots,u_{p,a_p}} \cup \set{u_{1,b_1},\dots,u_{p,b_p}}\right)
    $$
    hence $F_i \cap F_j \subseteq F \cap F_j$. 
    Finally, we have $|F \cap F_j| = \dim(\calS_{even}(T))$ by Lemma~\ref{lem. F1 cap F2 size} as $\nu_T(F_j)$ and $\nu_T(F)$ only differ by one digit.\newline

    \noindent\emph{Case 2}: $\ell(\nu_T(F_i)) = \ell(\nu_T(F_j))$ and $\nu_T(F_i) <_{lex} \nu_T(F_j)$.
    Since $\nu_T(F_i) <_{lex} \nu_T(F_j)$, there is an index $k \in [p]$ such that $a_k < b_k$ and $a_t = b_t$ for all $t < k$.
    Let 
    $$
    F = V_{even} \setminus \set{u_{1,b_1},\dots,u_{k-1,b_{k-1}}, u_{k,a_k}, u_{k+1,b_{k+1}},\dots,u_{p,b_p}} \in U_T 
    $$ 
    with
    $$
    \nu_T(F) = (b_1,\dots,b_{k-1},a_{k},b_{k+1},\dots,b_p).
    $$
    Then $F \in \mathcal{F}(\calS_{even}(T))$ by Lemma~\ref{lem. non-zero to 1 is TDS}.
    To show that $F <_T F_j$, we consider two cases.
    If $a_k = 1$, then $\ell(\nu_T(F)) = \ell(\nu_T(F_j)) + 1$, hence $F <_T F_j$.
    If $a_k > 1$, then $\ell(\nu_T(F)) = \ell(\nu_T(F_j))$ and $\nu_T(F) <_{lex} \nu_T(F_j)$ as $a_k < b_k$, hence $F <_T F_j$.
    Similar to Case 1, we get $F_i \cap F_j \subseteq F \cap F_j$, and $|F \cap F_j| = \dim(\calS_{even}(T))$ by Lemma~\ref{lem. F1 cap F2 size}.
\end{proof}

\end{document}
